\documentclass{article}
\begin{document}
\section{Introduction}
Thin liquid films evaporate quickly on heated walls. Their behaviour controls dryout in many compact heat exchangers. We still lack a clear picture of how wall wettability changes film rupture. In this work, we combine molecular dynamics with a simple film model to address this gap.

Earlier simulations focused on pool boiling \cite{a1}. Here, we look at thin films instead. The film thickness sets the time available for vaporisation. We therefore expect wettability to matter more for thin films than for pools.

\section{Results}
Figure 2 shows the film thickness against time for three contact angles. The film on the hydrophilic wall thins slowly and ruptures last. On the hydrophobic wall, a dry patch appears early. We find that rupture time falls as the contact angle grows.

We next look at the temperature field near the contact line. Figure 3 shows a hot spot under the thinnest part of the film. The hot spot is stronger on the hydrophobic wall. This agrees with the faster rupture seen in Figure 2.

The interfacial thermal resistance also changes with wettability. We measure it from the temperature jump at the wall. It is lower on the hydrophilic wall. Heat therefore crosses the interface more easily there, and the film stays stable for longer.

We now compare our results with the film model. The model captures the trend but underestimates the rupture time. The gap is largest for the most hydrophilic wall. We think the model misses the adsorbed layer that forms on this wall.

\section{Conclusions}
In this work, we studied thin-film evaporation on walls of different wettability. Rupture comes later on hydrophilic walls. The interfacial resistance is lower there, and heat crosses the wall more easily. These results may help in the design of wick structures.
\end{document}
